\scanpage{139}
\uline{Prop} \inserted{On suppose $2$ inv. sur $S$.}
Soit $H$ schéma en groupes lisse sur $S$,
de dim rel. $1$, et $X$ un torseur sous $H$.
Soit $(\unsure{\mathscr J|\mathscr V|\mathcal V},\unsure{\mathscr L|L|\mathcal L},\deleted{\unsure{O|\underline\omega}})$ l'Alg. de Lie de rang loc. $3$ canonique
sur $X$, \inserted{$(\mathbb P,\sigma)$ le fibré proj. à section associé,}
$H$ y opère par translation et le descend
par
\[
\underbrace{\underline\omega_0}_{\omega_0}\subset
\underbrace{L_0}_{\underline O_S}\subset
\underbrace{\mathscr J_0}_{\omega_0^{-1}}\ \text{sur } X,\qquad
(\mathbb P_0,\sigma_0)\quad \mathbb P_0=\mathbb P(L_0).
\]
\[
L_0\deleted{\unsure{\simeq|=}}\simeq(\underline O_H/\underline{\mathscr J}^3)^+,\quad
\underline{\mathscr J}\ \text{id. d'augm de }H
\]

Alors les conn. affines \inserted{sur $X$} invariantes sous $H$
sont en corr. $1$-$1$ avec \inserted{les sections d'}un certain sous
torseur \deleted{\unsure{sur}} de $\mathscr J_0\otimes\underline\omega_0$ sous $\underline\omega_0\otimes\underline\omega_0=\underline\omega_0^2$,
se projettant sur la section $1$ de $\underline\omega_0^{-1}\otimes\underline\omega_0\simeq\underline O_S$,
[formé des sections \inserted{locales} de $\mathscr J_0\otimes\underline\omega_0$ qui, pour
un \unsure{splittge|splitting} convenable de $\mathscr J_0$ \unsure{relmt} à $\underline\omega_0$, \inserted{(série de Hausdorff)} i.e.
\[
\mathscr J_0\simeq\underline\omega_0\oplus\underline O_S\oplus\underline\omega_0^{-1}
\]
\margin{splitting aussi donné par $\underline O_S\subset L_0$ correspond$^{\mathrm t}$ aux jets d'ordre $2$ de trivialisations \deleted{\unsure{canoniques}} ; ou ceux $\underline O_S\otimes\underline\omega$ \unsure{induis$^{\mathrm t}$} : $\omega_0\longrightarrow L_0\otimes\underline\omega_0$}
s'écrivent sous la forme
\[
\delta_2+0+1\quad\in\Gamma(\underline\omega_0^2+\underline\omega_0+\underline O_S)
\]

Les sections de ce sous-torseur correspondent
aux sections de $\mathbb P(\omega_0+\omega_0^{-1})$ qui sont disjointes
de la section canonique évidente de
$\omega_0\subset\omega_0+\omega_0^{-1}$, -- encore aux splittings
de $0\longrightarrow\omega_0\longrightarrow\omega_0+\omega_0^{-1}\longrightarrow\omega_0^{-1}\longrightarrow0$.
Il y en a une canonique (prenant pour celle-ci
celle de \unsure{Kampbell|Campbell}-Hausdorff) correspondant à la section $0+0+1$ de
$\mathscr J_0\otimes\underline\omega_0$. Elle est \uline{caractérisée} \inserted{parmi les conn. affines inv.} par la condition
que \deleted{\unsure{si $P$}}\inserted{loc. sur $S$ ($\beta$)} l'on puisse trouver sur
$P$ une \inserted{$G_a$-}réduction \inserted{invariante par translations} du groupe structural
à $G_a$, i.e. qu'on puisse trouver une
% margin note written sideways
\margin{$\beta$ Si on veut quelque chose de global sur $S$, il faut remplacer $G_a$ par $V(\underline\omega)$}
% archive sheet number: 138
